\section{INTRODUCTION}
\label{sec:intro}

Low-Power and Lossy Networks (LLNs) form the communication backbone of the Internet of Things (IoT), enabling applications from precision agriculture and environmental monitoring to industrial control and healthcare~\cite{rfc6550,rfc8200,rfc4944}. These networks operate under severe constraints: limited memory, strict energy budgets, low data rates, and intrinsically unreliable radio links. The IETF standardized the IPv6 Routing Protocol for Low-Power and Lossy Networks (RPL) to organize such networks into Destination-Oriented Directed Acyclic Graphs (DODAGs), allowing multiple objective functions (OFs) and routing metrics to reflect deployment-specific goals~\cite{rfc6550,rfc6551,gaddour2012rplsurvey}.

Despite its flexibility, a single-instance RPL deployment with a default OF collapses heterogeneous application requirements into one rank semantics. Emergency or control traffic experiences head-of-line blocking behind bulk telemetry; energy-rich nodes are routed identically to critically depleted ones; and the routing plane remains blind to link-reliability degradation from internal interference or early-stage faults~\cite{rfc7416}. These gaps are not theoretical: they arise whenever a single-path, single-metric routing core simultaneously serves safety-critical, best-effort, and long-lived maintenance flows.

Two complementary research directions have emerged to address these limitations. RPLMQoS~\cite{belacel2025rplmqos} belongs to the family of traffic-aware RPL extensions~\cite{nassar2018qosrpl,solapure2024rplofs,wakili2024mlrpl}: it adjusts path cost using composite link and end-to-end indicators (Expected Transmission Count---ETX, latency, jitter) and associates them with a traffic class so that rank computation reflects application urgency. In parallel, AER-RPL~\cite{belacel2025aerrpl} (Adaptive Energy-aware RPL) belongs to energy-responsive RPL designs~\cite{lamaazi2018energyrpl,limjoco2018analytical,wang2024drplcomnet,poornima2025fuzzyclrpl}, emphasizing residual energy, harvesting-aware predictions, and lightweight trust signals to keep routes viable under supply variability~\cite{farooq2022trustrpl,pishdar2021pccrpl}. Together with foundational contributions~\cite{bhandari2020multitopology}, these prior works establish the building blocks for a unified protocol.

Until now, QoS and energy awareness have been pursued in partial isolation: a QoS-centric stack may overlook energy-collapse dynamics, while an energy-centric stack may starve latency-sensitive flows. This paper proposes \textbf{AER-MQoS} (\emph{Adaptive Energy-aware Multi-Objective QoS Routing}), a disciplined unification that fuses traffic-class-aware QoS differentiation, energy-aware parent selection, and link-reliability-informed routing into a single coherent framework. The protocol retains RPL's control-plane structure~\cite{rfc6550,rfc9010} and introduces:

\begin{itemize}
  \item A \textbf{context fusion model} $(\alpha,\beta,\gamma)$ that makes the trade-off between QoS responsiveness ($\alpha$) and energy preservation ($\beta$) explicit. The bounded context coefficient $\gamma$ depends jointly on the application domain (health, industry, agriculture, environment) and the active traffic class among $\{C0,\ldots,C3\}$.
  \item A \textbf{Multi-Criteria Score (MCS)} that aggregates normalized ETX, an ETX-derived delay surrogate, normalized residual energy (NRE), predicted depletion risk, and a compact link-reliability trust signal into a scalar suitable for rank-based path selection.
  \item A \textbf{dedicated objective function} (experimental OCP~8) that couples the MCS to MRHOF-style hysteresis (Minimum Rank with Hysteresis Objective Function)~\cite{rfc6719} whenever \texttt{rpl-of-aer-plus.c} is linked.
  \item \textbf{Application-layer WRR queues} ($6{:}3{:}2{:}1$ for $C3{:}C2{:}C1{:}C0$) that reduce inter-class interference above UDP while keeping RAM costs explicit and bounded.
  \item A \textbf{lightweight Q-learning nudge} (Section~\ref{sec:ql}) that can adjust $\gamma$ within $\pm15$ hundredths based on recent transmission success and battery stress---documented for traceability and future ablation.
\end{itemize}

The protocol is evaluated through reproducible Cooja simulations with $N{=}4$ independent seeds, four firmware variants, and lossless/lossy UDGM channels. All variants exceed $93\%$ PDR (per-seed $93.4$--$99.8\%$) under lossy conditions; per-class PDR on seed~20260603 provides observational evidence that AER-MQoS differentiates traffic classes ($95.96\%$~C3 vs~$79.52\%$~C0) compared to MRHOF ($92.73\%$~C3 vs~$96.39\%$~C0), alongside episodic sensitivity on the same seed that motivates larger-scale confirmatory campaigns. The primary contribution is an open, inspectable, and extensible framework that unifies QoS, energy, and trust concerns in a single RPL objective function, backed by a transparent CSV-to-figure pipeline designed for community reuse and confirmatory extension.

\noindent\textbf{Article organization}\par
Section~\ref{sec:related} positions AER-MQoS within standardized RPL and representative extensions. Section~\ref{sec:arch} details the modular architecture, traffic classes, the MCS, and the OCP~8 objective function. Sections~\ref{sec:sec}--\ref{sec:ql} instantiate trust, energy, and the Q-learning specification. Section~\ref{sec:eval} reports the Cooja-only evaluation protocol, results, and reproducibility artifacts. Section~\ref{sec:concl} synthesizes contributions, limitations, and the future roadmap.